\section{Prompt 工程与交互策略}
%% ============================================================

如何与 Claude 对话直接决定输出质量。以下技巧帮助你写出更有效的 prompt。

\subsection{用 ultrathink 触发深度推理（Tip \#7）}

\texttt{ultrathink} 是一个关键词，会将 effort 设为最高并触发 Opus 4.6 的自适应推理模式。Claude 会根据问题复杂度动态分配 thinking tokens。

\begin{importantbox}{何时使用 ultrathink}
适用于：架构决策、复杂 debug、多步推理。不适用于：变量重命名等简单任务（浪费 thinking tokens）。也可以用 \texttt{/effort} 命令永久设置推理强度。
\end{importantbox}

\subsection{用 Plan Mode 规划复杂任务（Tip \#11）}

对于多文件变更、不熟悉的代码、架构决策，先进入 Plan Mode（\texttt{Shift+Tab} 切换）。虽然多花几分钟，但能防止 Claude 花 20 分钟自信地解决错误的问题。

对于一句话就能描述的小改动，直接执行即可，跳过规划。

\subsection{直接贴原始数据，不要自己解释 Bug（Tip \#13）}

\begin{warningbox}{常见误区：用文字描述 Bug}
用文字描述 bug 往往会丢失 Claude 定位根因所需的细节。直接粘贴 error log、CI 输出或 Slack 对话，然后说"fix"。Claude 擅长从原始日志中追踪问题。
\end{warningbox}

也可以直接从终端 pipe 输出：

\begin{lstlisting}[caption={管道传输错误信息}, language=bash]
npm test 2>&1 | claude "fix the failing tests"
\end{lstlisting}

\subsection{用 @ 指定文件路径（Tip \#25）}

使用 \texttt{@src/auth/middleware.ts} 直接引用文件。Claude 虽然能自己搜索代码库，但每一步搜索都消耗 tokens 和上下文。直接指定文件跳过整个搜索过程。

\subsection{用模糊 Prompt 探索不熟悉的代码（Tip \#26）}

"What would you improve in this file?" 是一个很好的探索型 prompt。不是每个 prompt 都需要精确。当你想要新视角审视现有代码时，模糊问题给 Claude 发挥空间，能发现你自己不会想到去问的问题。

\subsection{让 Claude 反过来采访你（Tip \#44）}

当你无法完整描述需求时，让 Claude 来提问：

\begin{importantbox}{Interview 模式}
告诉 Claude："I want to build [feature]. Interview me to understand the full requirements before writing any code. Ask me one question at a time." Claude 会逐步引导你补全需求细节。完成后，开一个新会话用干净上下文执行。
\end{importantbox}

\subsection{用语音输入写更丰富的 Prompt（Tip \#23）}

运行 \texttt{/voice} 启用按键说话功能，按住空格键口述。语音 prompt 自然比打字包含更多上下文，因为口述时你会不自觉地解释背景、提到约束条件。需要 Claude.ai 账号（非 API key）。可以在 \texttt{\textasciitilde/.claude/keybindings.json} 中自定义按键。

\subsection{本章小结}

Prompt 工程的核心原则：复杂任务用 \texttt{ultrathink} 和 Plan Mode；贴原始数据而非文字描述；用 \texttt{@} 精确指向文件；在需求不明确时让 Claude 反向采访你。匹配推理强度与任务复杂度是控制成本的关键。

%% ============================================================
